\Problem{Imprimir expressions}

\Statement

En aquest exercici considerarem arbres que representen expressions 
sobre els operadors \verb|+|, \verb|-|, \verb|*|, i sobre operands 
que són nombres naturals. Per exemple, el següent arbre 

\begin{verbatim}
*
|-- -
|   |-- 7
|   |-- 1
|   |-- 5
|   '-- 2
'-- +
    |-- 7
    |-- 3
    '-- 8
\end{verbatim}

representa l'expressió \verb|((7 - 1 - 5 - 2) * (7 + 3 + 8))|.

\medskip

Tots els nodes de l'arbre que són fulles contenen un valor natural,
i els nodes amb fills sempre en tenen almenys 2 i el seu valor és 
sempre un operador. No hi ha subarbres buits.

Implementeu, doncs, la funció següent:

\begin{verbatim}
/**
 * @brief Transforma una expressió en la seva representació com a `string`.
 *
 * Una expressió està formada per operadors (`+`, `*` i `-`), i
 * operands (naturals), que són nodes d'un arbre. Els operadors tenen dos
 * fills o més, i els operands són fulles (no tenen fills).
 *
 * L'expressió representada com a `string` és de la següent manera. Per a
 * operands: cal retornar l'operand mateix. Per a operadors, cal retornar
 * els fills de l'operador, separats per l'operador, amb un espai entre
 * operands i operadors, i tot el conjunt sempre entre parèntesis.
 *
 * @pre L'arbre representa una expressió ben formada
 *
 * @param t L'arbre que representa l'expressió.
 * @returns La representació com a `string` de `t`
 */
string expression_to_string(Tree<string> t);  
\end{verbatim}

\Observation

Els fitxers públics (icona del gatet) contenen:

\medskip

\begin{tabular}{ll}
 \verb|tree.hh| & la classe \verb|Tree| \\
 \verb|tree-io.hh| & l'entrada/sortida de \verb|Tree| \\
 \verb|main.cc| & el programa principal \\
\end{tabular}

\medskip

També hi ha un \verb|Makefile| i el directori \verb|.vscode| 
que té la configuració per compilar i depurar amb VSCode.

\medskip

Cal implementar \verb|expression_to_string| en un 
\textbf{fitxer \texttt{.cc} nou}, compilar, i finalment 
\textbf{enviar només el fitxer amb la funció}.

\Input

Cada cas consisteix en una representació textual d'una expressió del 
tipus definit. (Aquesta lectura ja la fa el programa principal.)

\Output

Per a cada cas, la sortida és el resultat de transformar l'expressió 
a la seva representació textual, amb cada resultat en una línia separada. 
(Això també ho fa el programa principal.)

\SampleTwoCol

